\documentclass[14pt,a4paper]{extarticle}
\usepackage{bsuir_stp}

\ifxetexorluatex
  \newfontfamily\stpcodefont{Courier New}
  \newcommand{\straightquote}{\symbol{34}}
\else
  \newcommand{\straightquote}{{\fontencoding{T1}\selectfont\symbol{34}}}
  \usepackage[T1,T2A]{fontenc}
  \newcommand{\stpcodefont}{\fontencoding{T1}\fontfamily{pcr}\selectfont}
  \renewcommand{\texttt}[1]{{\stpcodefont #1}}
\fi
\lstset{
  language=Python,
  basicstyle=\stpcodefont\fontsize{13pt}{15.6pt}\selectfont,
  keywordstyle=\bfseries,
  columns=fullflexible,
  aboveskip=6pt,
  belowskip=6pt,
  xleftmargin=0pt
}
\graphicspath{{images/}}

\stpFaculty{Факультет компьютерных систем и сетей}
\stpDepartment{Кафедра ПОИТ}
\stpSubject{Методы машинного обучения}
\stpLabNumber{1}
\stpLabTopic{Основы Python}
\newcommand{\variant}{3}
\stpVariant{\variant}
\stpStudent{Беседин А.А.}
\stpGroup{351002}
\stpTeacher{Лебедевич А.В.}
\stpCity{Минск}
\stpYear{2026}

\newcommand{\shot}[3]{%
  \begin{figure}[H]
    \centering
    \includegraphics[width=#3\textwidth,height=0.8\textheight,keepaspectratio]{#1}
    \caption{#2}
    \label{fig:#1}
  \end{figure}%
}

\begin{document}

\makestptitlepage

\section{Цель работы}
Целью лабораторной работы является освоение базовых возможностей языка программирования Python, применяемого в дальнейшем для решения задач машинного обучения: изучение встроенных типов данных и коллекций, управляющих конструкций, функций и механизма обработки исключений, а также принципов объектно-ориентированного программирования (инкапсуляции, наследования и полиморфизма). Практическая часть работы заключается в разработке и тестировании набора консольных программ, решающих задачи индивидуального задания.

\section{Краткие теоретические сведения}

\subsection{Язык Python и базовые типы данных}
Python -- высокоуровневый интерпретируемый язык программирования общего назначения с динамической строгой типизацией и автоматическим управлением памятью. Благодаря простоте синтаксиса и развитой экосистеме библиотек (NumPy, pandas, scikit-learn, PyTorch) Python является основным языком в области анализа данных и машинного обучения.

Любое значение в Python является объектом, а переменная -- лишь ссылкой на объект. Тип переменной определяется во время выполнения программы по значению, которое ей присвоено. К основным встроенным типам относятся целые числа \texttt{int} (произвольной точности), вещественные числа \texttt{float}, логический тип \texttt{bool} со значениями \texttt{True} и \texttt{False} и строки \texttt{str}.

Строка -- неизменяемая упорядоченная последовательность символов Юникода. Строки поддерживают индексацию \texttt{s[i]}, срезы \texttt{s[start:stop:step]} (например, \texttt{s[::-1]} возвращает перевернутую строку) и большое число методов. В работе используются методы проверки символов: \texttt{isdigit()} (цифра), \texttt{isalpha()} (буква), \texttt{isupper()} и \texttt{islower()} (прописная и строчная буква), а также \texttt{upper()} и \texttt{split()}.

Ввод данных с клавиатуры выполняется функцией \texttt{input()}, которая всегда возвращает строку, поэтому для получения чисел используется явное преобразование типов \texttt{int()} и \texttt{float()}. Вывод выполняется функцией \texttt{print()}. Функция \texttt{map()} применяет заданное преобразование к каждому элементу последовательности, что позволяет компактно прочитать ряд чисел: \texttt{list(map(int, input().split()))}.

\subsection{Коллекции}
Список \texttt{list} -- изменяемая упорядоченная коллекция элементов произвольных типов. Списки поддерживают добавление элементов (\texttt{append}), подсчет вхождений (\texttt{count}), поиск индекса первого вхождения (\texttt{index}, в том числе с заданной начальной позицией) и срезы.

Кортеж \texttt{tuple} -- неизменяемая упорядоченная коллекция. Кортежи применяются для хранения фиксированных наборов значений и для возврата нескольких результатов из функции: запись \texttt{return even, odd} фактически возвращает кортеж, который затем распаковывается в несколько переменных.

Словарь \texttt{dict} -- изменяемое отображение уникальных хешируемых ключей на значения, реализованное в виде хеш-таблицы. Доступ к значению по ключу выполняется в среднем за время $O(1)$. Обход словаря в цикле \texttt{for} перебирает ключи, метод \texttt{items()} возвращает пары «ключ -- значение», а оператор \texttt{in} проверяет наличие ключа. Начиная с версии 3.7 словари сохраняют порядок вставки элементов.

\subsection{Управляющие конструкции и функции}
Ветвление реализуется конструкцией \texttt{if}~--~\texttt{elif}~--~\texttt{else}. Для организации циклов служат оператор \texttt{for}, перебирающий элементы любого итерируемого объекта (в том числе диапазона \texttt{range(n)}), и оператор \texttt{while}, выполняющий тело, пока истинно условие. Оператор \texttt{break} досрочно завершает цикл, \texttt{continue} -- переходит к следующей итерации.

Функция объявляется ключевым словом \texttt{def}, принимает параметры и возвращает результат оператором \texttt{return}. Функции позволяют выделить логически завершенные фрагменты алгоритма и использовать их повторно. Встроенная функция \texttt{isinstance(obj, cls)} проверяет принадлежность объекта к типу, что позволяет реализовать функцию, поведение которой зависит от типа аргумента.

Примером алгоритма, использующего цикл \texttt{while}, является проверка числа на простоту. Натуральное число $n \geq 2$ простое, если оно не имеет делителей $d$, удовлетворяющих условию
\begin{equation}
  2 \leq d \leq \sqrt{n},
\end{equation}
поэтому перебор делителей достаточно вести до $\sqrt{n}$, что дает сложность алгоритма $O(\sqrt{n})$.

\subsection{Обработка исключений}
Исключение -- объект, сигнализирующий об ошибке во время выполнения программы. Если исключение не обработано, выполнение программы прерывается. Для перехвата исключений используется конструкция \texttt{try}~--~\texttt{except}~--~\texttt{else}~--~\texttt{finally}: в блоке \texttt{try} размещается потенциально опасный код, блоки \texttt{except} перехватывают исключения указанных типов, блок \texttt{else} выполняется при отсутствии ошибок, а блок \texttt{finally} выполняется всегда, независимо от того, возникло ли исключение. Типичными встроенными исключениями являются \texttt{ValueError} (некорректное значение, например, при преобразовании строки «abc» в число) и \texttt{ZeroDivisionError} (деление на ноль).

\subsection{Объектно-ориентированное программирование}
Класс -- пользовательский тип данных, описывающий состояние (атрибуты) и поведение (методы) объектов. Объект (экземпляр класса) создается вызовом класса, при этом автоматически вызывается метод-инициализатор \texttt{\_\_init\_\_}, в котором задаются атрибуты экземпляра. Первый параметр метода \texttt{self} -- ссылка на текущий объект.

Инкапсуляция -- объединение данных и методов в классе и ограничение прямого доступа к внутреннему состоянию. В Python она реализуется соглашениями об именовании: одно подчеркивание (\texttt{\_area}) обозначает защищенный атрибут, два подчеркивания (\texttt{\_\_make\_deal}) -- приватный, имя которого искажается интерпретатором (name mangling) и недоступно извне класса напрямую.

Наследование позволяет создать производный класс на основе базового, переняв его атрибуты и методы. Вызов \texttt{super().\_\_init\_\_()} передает управление инициализатору родительского класса. Полиморфизм -- возможность единообразно работать с объектами разных классов через общий интерфейс: производные классы переопределяют метод базового класса, и при вызове выполняется реализация, соответствующая фактическому типу объекта.

Атрибуты класса, объявленные в теле класса, разделяются всеми экземплярами. Методы класса, помеченные декоратором \texttt{@classmethod}, получают первым параметром сам класс \texttt{cls} и применяются для работы с атрибутами класса. Статические методы (\texttt{@staticmethod}) не получают ни экземпляра, ни класса и представляют собой обычные функции, логически сгруппированные внутри класса.

В качестве примера предметной логики класса рассмотрим треугольник со сторонами $a$, $b$, $c$. Треугольник существует, если выполняются неравенства
\begin{equation}
  a + b > c, \quad a + c > b, \quad b + c > a,
\end{equation}
а его площадь вычисляется по формуле Герона
\begin{equation}
  S = \sqrt{p\,(p - a)(p - b)(p - c)}, \quad p = \frac{a + b + c}{2},
\end{equation}
\noindent где $S$ -- площадь треугольника;\\
$p$ -- полупериметр треугольника;\\
$a$, $b$, $c$ -- длины сторон треугольника.

\section{Формулировка индивидуального задания}
В соответствии с вариантом №\,\variant{} требуется разработать на языке Python набор консольных программ, состоящий из трех частей.

\subsection{Часть 1. Базовые конструкции и коллекции}
\begin{enumerate}[label=\arabic*.]
  \item Определить, сколько в числе четных цифр, а сколько нечетных. Число вводится с клавиатуры. – 2 балла
  \item Посчитать, сколько пар (стоят рядом) верхнего и нижнего регистра находится в веденном с клавиатуры слове. (Пример HjkLM- 1 пара нижнего, 1 пара верхнего), а также сколько всего букв в слове. – 2 балла
  \item Найдите сумму отрицательных элементов списка. Найдите сумму элементов списка между двумя первыми нулями. Если двух нулей нет в списке, то выведите ноль. – 1 балл
  \item В настольной игре Скрабл (Scrabble) каждая буква имеет определенную ценность. В случае с английским алфавитом очки распределяются так: A, E, I, O, U, L, N, S, T, R – 1 очко; D, G – 2 очка; B, C, M, P – 3 очка; F, H, V, W, Y – 4 очка; K – 5 очков; J, X – 6 очков; Q, Z – 10 очков.

  Напишите программу, которая вычисляет стоимость введенного пользователем слова. – 2 балла
  \item Реализуйте программу «Магазин автозапчастей», которая будет включать в себя шесть пунктов меню. У вас есть словарь, где ключ – название продукции. Значение – список, который содержит состав продукции, цену и кол-во (шт),которое есть в магазине. 1. Просмотр описания: название – описание 2. Просмотр цены: название – цена. 3. Просмотр количества: название – количество. 4. Всю информацию. 5. Покупка В пункте «Покупка» необходимо совершить покупку, с клавиатуры вводите название продукции и его кол-во, n – выход из программы. Посчитать цену выбранных товаров и сколько товаров осталось в изначальном списке 6. До свидания – 2 балла
  \item В кортеже целых чисел найдите максимальный и минимальный элементы. – 1 балл
\end{enumerate}

\subsection{Часть 2. Функции и обработка исключений}
\begin{enumerate}[label=\arabic*.]
  \item Напишите функцию is\_password\_good(password), которая принимает в качестве аргумента строковое значение пароля password и возвращает значение True если пароль является надежным и False в противном случае. Пароль является надежным если: – его длина не менее 8 символов; – он содержит как минимум одну заглавную букву (верхний регистр); – он содержит хотя бы одну цифру. – 2 балла
  \item Напишите функцию, которая будет принимать один аргумент. Если в функцию передаётся словарь, то вывести ключ с максимальным значением. Если список, то найти количество элементов, расположенных до первого отрицательного элемента. Число – проверить простое, или нет. Строка – вывести в обратном порядке. Найти сумму цифр в строке. Сделать проверку со всеми этими случаями. – 4 балла
  \item Даны два числа n и m. Создайте двумерный массив размером n$\times$m и заполните его символами \straightquote.\straightquote{} и \straightquote*\straightquote{} в шахматном порядке. В левом верхнем углу должна стоять точка. – 2 балла
  \item Напишите программу, демонстрирующую работу try\textbackslash{}except\textbackslash{}finally – 2 балла
\end{enumerate}

\subsection{Часть 3. Объектно-ориентированное программирование}
Задание 1: Класс Triangle Создайте класс Triangle. В нём пропишите 3 (метода) функции. Первый метод: проверка на существование треугольника по данным сторонам. Второй метод : нахождение площади треугольника. Третий метод: нахождение периметра треугольника. – 1 балл

Задание 2: Класс House 1. Создайте класс House 2. Создайте метод \_\_init\_\_() и определите внутри него два динамических свойства: \_area и \_price. Свои начальные значения они получают из параметров метода \_\_init\_\_() 3. Создайте метод final\_price(), который принимает в качестве параметра размер скидки и возвращает цену с учетом данной скидки. Класс SmallHouse 1. Создайте класс SmallHouse, унаследовав его функционал от класса House 2. Внутри класса SmallHouse переопределите метод \_\_init\_\_() так, чтобы он создавал объект с площадью 40м2 Класс Human 1. Реализуйте приватный метод make\_deal(), который будет отвечать за техническую реализацию покупки дома: уменьшать количество денег на счету и присваивать ссылку на только что купленный дом. В качестве аргументов данный метод принимает объект дома и его цену. 2. Реализуйте метод buy\_house(), который будет проверять, что у человека достаточно денег для покупки, и совершать сделку. Если денег слишком мало - нужно вывести предупреждение в консоль. Параметры метода: ссылка на дом и размер скидки – 3 балла

Задание 3: Создать классы Circle (круг), Square (квадрат), Rectangle (прямоугольник) для описания плоских геометрических фигур. Реализовать метод нахождения площади фигуры. Переопределить метод нахождения площади фигуры.

Задание 4:\newline
Придумать класс самостоятельно, реализовать в нем методы экземпляра класса, статические, методы, методы класса. – 1 – 3 балла

\section{Код решения индивидуального задания}
Ниже приведен полный исходный код программ, разработанных для решения индивидуального задания.

\subsection{Часть 1. Базовые конструкции и коллекции}
\subsubsection{Задание 1}
\lstinputlisting{code/sublab1_task1.py}
\subsubsection{Задание 2}
\lstinputlisting{code/sublab1_task2.py}
\subsubsection{Задание 3}
\lstinputlisting{code/sublab1_task3.py}
\subsubsection{Задание 4}
\lstinputlisting{code/sublab1_task4.py}
\subsubsection{Задание 5}
\lstinputlisting{code/sublab1_task5.py}
\subsubsection{Задание 6}
\lstinputlisting{code/sublab1_task6.py}

\subsection{Часть 2. Функции и обработка исключений}
\subsubsection{Задание 1}
\lstinputlisting{code/sublab2_task1.py}
\subsubsection{Задание 2}
\lstinputlisting{code/sublab2_task2.py}
\subsubsection{Задание 3}
\lstinputlisting{code/sublab2_task3.py}
\subsubsection{Задание 4}
\lstinputlisting{code/sublab2_task4.py}

\subsection{Часть 3. Объектно-ориентированное программирование}
\subsubsection{Задание 1}
\lstinputlisting{code/sublab3_task1.py}
\subsubsection{Задание 2}
\lstinputlisting{code/sublab3_task2.py}
\subsubsection{Задание 3}
\lstinputlisting{code/sublab3_task3.py}
\subsubsection{Задание 4}
\lstinputlisting{code/sublab3_task4.py}

\section{Результаты выполнения индивидуального задания}
Каждая программа была запущена из командной строки на нескольких наборах входных данных, включающих как типичные, так и граничные случаи.

\subsection{Часть 1. Базовые конструкции и коллекции}
Результат работы программы подсчета четных и нечетных цифр приведен на рисунке~\ref{fig:s1t1}. Для числа 4815162342 найдено 6 четных и 4 нечетные цифры, для числа 13579 четные цифры отсутствуют.

\shot{s1t1}{Подсчет четных и нечетных цифр числа}{0.85}

На рисунке~\ref{fig:s1t2} показан подсчет пар соседних строчных и прописных букв и общего количества букв в слове.

\shot{s1t2}{Подсчет пар строчных и прописных букв в слове}{0.85}

Вычисление суммы отрицательных элементов и суммы элементов между первыми двумя нулями приведено на рисунке~\ref{fig:s1t3}. Во втором запуске нулевые элементы отсутствуют, поэтому вторая сумма равна нулю.

\shot{s1t3}{Вычисление сумм элементов списка}{0.85}

На рисунке~\ref{fig:s1t4} приведен подсчет стоимости слов: слово Quiz оценено в $10 + 1 + 1 + 10 = 22$ очка.

\shot{s1t4}{Подсчет стоимости слова по правилам игры «Эрудит»}{0.85}

Работа пунктов меню программы «Магазин автозапчастей», отвечающих за просмотр информации о товарах, а также реакция на некорректный пункт меню показаны на рисунке~\ref{fig:s1t5a}.

\shot{s1t5a}{Просмотр информации о товарах в программе «Магазин автозапчастей»}{0.75}

Процесс покупки товаров приведен на рисунке~\ref{fig:s1t5b}. Программа отклоняет покупку при недостаточном количестве товара на складе и при вводе несуществующего названия, подсчитывает итоговую сумму покупки ($2 \cdot 2300 + 2 \cdot 7800 = 20200$) и выводит обновленные остатки.

\shot{s1t5b}{Покупка товаров в программе «Магазин автозапчастей»}{0.75}

Поиск максимального и минимального элементов кортежа показан на рисунке~\ref{fig:s1t6}.

\shot{s1t6}{Поиск максимального и минимального элементов кортежа}{0.85}

\subsection{Часть 2. Функции и обработка исключений}
Проверка надежности паролей приведена на рисунке~\ref{fig:s2t1}: пароль qwerty слишком короткий, пароль password123 не содержит прописных букв, и только пароль Password123 удовлетворяет всем требованиям.

\shot{s2t1}{Проверка надежности паролей}{0.85}

На рисунке~\ref{fig:s2t2} показана обработка аргументов различных типов одной функцией: словаря, списка, целого числа и строки.

\shot{s2t2}{Обработка аргументов различных типов}{1.0}

Вывод шахматных досок размером $5 \times 8$ и $3 \times 3$ приведен на рисунке~\ref{fig:s2t3}.

\shot{s2t3}{Формирование шахматной доски}{0.85}

Обработка исключений при делении чисел показана на рисунке~\ref{fig:s2t4}: корректное деление, деление на ноль и ввод нечислового значения. Во всех трех случаях выполняется блок \texttt{finally}.

\shot{s2t4}{Обработка исключений при делении чисел}{0.85}

\subsection{Часть 3. Объектно-ориентированное программирование}
На рисунке~\ref{fig:s3t1} приведена работа класса \texttt{Triangle}: для сторон 3, 4, 5 вычислены площадь и периметр, а для сторон 1, 2, 10 определено, что треугольник не существует.

\shot{s3t1}{Работа класса Triangle}{0.85}

Покупка дома объектом класса \texttt{Human} показана на рисунке~\ref{fig:s3t2}. В первом случае денег достаточно для покупки дома с учетом скидки 10~\%, во втором -- покупка отклоняется.

\shot{s3t2}{Покупка дома объектом класса Human}{0.85}

Вычисление площадей фигур через единый метод \texttt{area()} базового класса приведено на рисунке~\ref{fig:s3t3}.

\shot{s3t3}{Демонстрация полиморфизма при вычислении площадей фигур}{0.85}

Работа класса \texttt{BankAccount} показана на рисунке~\ref{fig:s3t4}. В первом запуске после пополнения на 500 и снятия 200 на баланс 1300 начислено 10~\%, во втором -- статический метод отклоняет отрицательную сумму пополнения, а снятие суммы, превышающей баланс, запрещается.

\shot{s3t4}{Работа класса BankAccount}{0.85}

\section{Выводы}
В ходе выполнения лабораторной работы №\,1 были получены следующие результаты:
\begin{enumerate}
  \item Изучены встроенные типы данных и коллекции языка Python (строки, списки, кортежи, словари), а также основные управляющие конструкции. На их основе разработаны программы обработки чисел, строк и списков и интерактивная программа «Магазин автозапчастей» с меню.
  \item Освоены объявление и использование функций, возврат нескольких значений и проверка типа аргумента с помощью \texttt{isinstance}. Изучен механизм обработки исключений \texttt{try}~--~\texttt{except}~--~\texttt{finally}, позволяющий программе корректно реагировать на ошибочный ввод и деление на ноль.
  \item Изучены принципы объектно-ориентированного программирования: разработаны классы, демонстрирующие инкапсуляцию (защищенные и приватные члены), наследование, полиморфизм, а также использование атрибутов класса, методов класса и статических методов.
  \item Все программы протестированы на типичных и граничных наборах входных данных, результаты их работы совпадают с ожидаемыми. Полученные навыки программирования на Python будут использованы в последующих лабораторных работах по методам машинного обучения.
\end{enumerate}

\end{document}
